
%%%%%%%%%%%%%%%%%%%%%%%%%%%%%%%%%%%%%%%%%%%%%%%%%%%%%%%%%%%%%%
%%%%%% MA-0472 Introducción a la Geometría Diferencial %%%%%%%
%%%%%%%%%%%%%%%%%%%%%%%%%%%%%%%%%%%%%%%%%%%%%%%%%%%%%%%%%%%%%%
\newpage
\subsection*{MA-0472 Introducción a la Geometría Diferencial}
\addcontentsline{toc}{subsection}{MA-0472 Introducción a la Geometría Diferencial}

\begin{refsection}
\begin{table}[H]
\centering{
\begin{tabular}{|ll|}
\hline
%\multicolumn{2}{|c|}{\textbf{Nivel de virtualidad del curso:} Virtual}\\ \hline
\textbf{Año:} II  & \textbf{Requisitos:} MA-0351, MA-0261 y MA-0252  \\ 
\textbf{Ciclo:} IV  & \textbf{Correquisitos:} Ninguno \\ 
\textbf{Tipo de curso:} Teórico & \textbf{Horas presenciales por semana:}   \\ 
\textbf{Créditos:} 4 & \textbf{Horas de trabajo independiente por semana:}  \\ \hline
\end{tabular}
}
\end{table}



\subsubsection*{Descripción del curso}

El curso de curvas y superficies es el primero de dos cursos que  estudian los principios b\'asicos de la geometr\'{i}a diferencial, tanto cl\'asica como moderna.  

El estudio de curvas tiene dos aspectos que debemos considerar: el primero es el estudio local de ellas en el plano y en el espacio, donde  entre otros temas, se considera el aparato de Serret-Frenet y algunas de sus consecuencias como el teorema de Lancret. Y como un segundo aspecto se estudia, al menos una propiedad global: la llamada desigualdad isoperim\'etrica.  Cabe se\~analar que en esta parte se nota la dependencia de la teor\'{i}a de curvas en las ecuaciones diferenciales ordinarias.

El estudio de superficies, el tema principal del curso, es mucho m\'as profundo que el de curvas, tanto a nivel filos\'ofico como a nivel matem\'atico. El enfoque seguido en el curso es el Riemanniano, el que propuso Bernhard Riemann es su lectura inaugural en la Universidad de Gotingen: {\bf On the Hypothesis which Lie at the Foundations og Geometry  }. 



%Este es el primero de una secuencia de dos cursos que introducen los principios básicos del análisis matemático en una y varias variables reales. A lo largo de esta secuencia se cubren los temas usuales del análisis matemático
de la recta real y el cálculo diferencial e integral en una y varias variables. 

El curso se encuentra enfocado al desarrollo del formalismo y la abstracción de elementos del \'algebra lineal,  el c\'alculo en varias variables, y las ecuaciones diferenciales en la geometr\'{i}a, mediante la utilización de técnicas de razonamiento y demostración propias del área.

Se toma como punto de partida el estudio de la propiedad de completitud del conjunto de los números reales y algunas de sus implicaciones, además de las intuiciones y destrezas operacionales desarrolladas en los cursos MA-0251 Introducción al cálculo en una variable y MA-0351 Introducción al cálculo en varias variables, así como las bases de lógica y estrategias de demostración adquiridas en los cursos MA-0152 Matemática Exploratoria y MA-0252 Introducción a las demostraciones.





\subsubsection*{Objetivos}

\textbf{General:} Aplicar los conceptos básicos de completitud, convergencia y diferenciación para la demostración de propiedades del análisis en una variable real. \\

\textbf{Específicos:}

\begin{enumerate}
\item Asociar las propiedades derivadas de la completitud en los números reales con las respectivas nociones geométricas e intuitivas.
\item Demostrar resultados sobre convergencia de sucesiones y series.
\item Demostrar resultados sobre límites, continuidad, convergencia y diferenciabilidad de funciones de una variable real.
\item Contrastar los diferentes tipos de convergencia para sucesiones y series.
\item Aplicar los conceptos y resultados sobre completitud y convergencia en la construcción de las funciones trascendentes elementales.
\item Indagar en la literatura sobre temas que permitan complementar su formación.
\item Comunicar ideas matemáticas de manera clara y efectiva en forma oral y escrita.
\end{enumerate}

\subsubsection*{Contenidos}

\begin{enumerate}
\item \textbf{Completitud y propiedades de los números reales (2 semanas):} 
\begin{enumerate}
    \item Incompletitud de los números racionales.
    \item Completitud de los números reales: conjuntos acotados y axiomas del extremo superior e inferior.
    \item Aplicaciones: Principio de arquimedianidad, densidad de $\Q$ en $\R$, existencia de raíces $n$-ésimas, la no contabilidad de $\R$.
    \item Intervalos, valor absoluto, vecindarios y cercanía. 
\end{enumerate}

\item \textbf{Sucesiones y series numéricas (4 semanas):} 
\begin{enumerate}
    \item Nociones básicas, convergencia, álgebra de límites, sucesiones monótonas, recurrencia.
    \item Subsucesiones, limsup y liminf, Teorema de la Convergencia Monótona, el Teorema de los Intervalos Encajados, el Teorema de Bolzano-Weierstrass, sucesiones de Cauchy, puntos límite.
    \item Definición formal y convergencia de series.
    \item Demostración de los criterios de convergencia básicos.
    \item Los criterios de convergencia de Dirichlet y de Raabe.
    \item Convergencia condicional y absoluta. Reordenamientos de series.
    \item Producto de series.
\end{enumerate}

\item \textbf{Límites y continuidad (3 semanas):} 
\begin{enumerate}
    \item Límites de funciones. El criterio secuencial para la existencia de un límite. Propiedades básicas de los límites.
    \item Límites laterales, límites al infinito y límites infinitos. El Teorema de la Compresión o del Encaje.
    \item Definición de continuidad. Caracterización secuencial de la continuidad. Operaciones algebraicas con funciones continuas, composición de funciones continuas.
    \item Tipos de discontinuidades. Ejemplos de funciones continuas y con discontinuidades.
    \item Funciones continuas en intervalos cerrados y acotados. El Teorema de los Valores Extremos.
    \item Definición de continuidad uniforme. Criterio secuencial para la ausencia de continuidad uniforme. Continuidad uniforme en intervalos cerrados y acotados.
    \item El Teorema de los Valores Intermedios.
    \item Continuidad de la función inversa.
\end{enumerate}

\item \textbf{Diferenciación (2 semanas):} 
\begin{enumerate}
    \item Demostración de la Regla de la Cadena y del Teorema de la Función Inversa.
    \item Los Teoremas del Valor Medio: el Teorema de Rolle, el Teorema del Valor Medio, el Teorema del Valor Medio Generalizado. Consecuencias básicas de estos teoremas.
    \item El Teorema del Valor Extremo Interior, el Teorema de Darboux.
    \item La Regla de L'H\^{o}pital y sus variantes.
\end{enumerate}

\item \textbf{Sucesiones y series de funciones (4 semanas):} 
\begin{enumerate}
    \item Convergencia puntual y uniforme de sucesiones de funciones.
    \item La norma uniforme, el criterio de Cauchy para convergencia uniforme.
    \item Continuidad y derivabilidad de la función límite de una sucesión de funciones.
    \item Definición de convergencia puntual y uniforme de series de funciones.
    \item Continuidad y diferenciabilidad término a término de series de funciones.
    \item El Criterio de Cauchy para convergencia uniforme de series de funciones y el $M$-test de Weierstrass.
    \item Series de potencias. Radio e intervalo de convergencia, convergencia uniforme, el Teorema de Abel, diferenciabilidad, unicidad y Series de Taylor.
\end{enumerate}

\item \textbf{Construcción de las funciones trascendentes elementales (1 semana):} 
\begin{enumerate}
    \item Construcción de la función exponencial y logarítmica.
    \item Construcción de las funciones trigonométricas.
\end{enumerate}

\end{enumerate}

\subsubsection*{Metodología}

\noindent \textbf{Lineamientos metodológicos:} La persona docente a cargo de este curso debe respetar las siguientes pautas:
\begin{enumerate}
    \item La presentación de los contenidos debe priorizar la formalización de los conceptos estudiados en los cursos MA-0251 Introducción al cálculo en una variable y MA-0351 Introducción al cálculo en varias variables. En particular en este curso se debe asumir que las personas estudiantes  ya manejan las intuiciones y las habilidades operacionales sobre los conceptos básicos del cálculo diferencial e integral.
    \item En este curso se aplican las bases de lógica y las estrategias de demostración vistas en el curso MA-0252 Introducción a las demostraciones.
    \item Las demostraciones deben presentarse tomando en cuenta que están dirigidas a una persona estudiante principiante. En particular, se debe priorizar dar un mayor nivel de detalle por encima de la brevedad y la estética, cuando esto ayude a mejorar la comprensión de los temas.
    \item Cuando la naturaleza del contenido lo permita, se deben utilizar representaciones visuales que apoyen la comprensión de los temas.
\end{enumerate}

\textbf{Sugerencias metodológicas:} Adicionalmente, se tienen las siguientes recomendaciones para la persona docente a cargo de este curso:
\begin{enumerate}
    \item Utilizar el recurso tecnológico para reforzar distintos conceptos estudiados en el curso por medio de la visualización, cálculo y experimentación. Se pone a disposición de la persona docente un repositorio de herramientas listas para ser implementadas en el curso, tanto en las clases como para el trabajo independiente de las personas estudiantes.
    \item Aprovechar momentos adecuados del curso para motivar el estudio por medio de reseñas acerca del desarrollo histórico del análisis matemático y su importancia en aplicaciones dentro y fuera de la matemática
    \item En la evaluación, se sugiere utilizar estrategias que promuevan el trabajo constante de parte de las personas estudiantes a lo largo del ciclo lectivo. Por ejemplo, esto se puede hacer por medio de tareas o quizzes.
    \item Para fomentar el desarrollo de las habilidades investigativas en el estudiantado, se sugiere la implementación de pequeños proyectos de investigación. Por ejemplo, pueden ser proyectos de investigación bibliográfica o también proyectos cortos donde se desarrolle un tema por medio de ejercicios dirigidos, como se propone, por ejemplo, en el libro de Abbott \cite{Abb15}.
    \item Dado que hoy en día la mayoría del trabajo en matemática, tanto a nivel académico como profesional, es de naturaleza colaborativa, se sugiere a la persona docente implementar estrategias que promuevan el trabajo en equipo.
\end{enumerate}



%%%%%%%%%%%%%%%%%%%%%%%%%%%%%%%%%%
%%%%%%%%%%%%%%%%%%%%%%%%%%%%%%%%%%
%% No cite para las referencias %%
\nocite{doC16}
\nocite{Kuh15}
%%%%%%%%%%%%%%%%%%%%%%%%%%%%%%%%%%
%%%%%%%%%%%%%%%%%%%%%%%%%%%%%%%%%%


\printbibliography[heading=subbibliography]
\end{refsection}
